\chapter{Минимальный локальный перенос и запрет массового замыкания}

Предыдущая глава заменила поиск очередного статического потенциала задачей
вывода локального оператора переноса. Настоящий гейт выполняет первый
вычислимый тест этой программы. Проверяется наиболее благоприятный
минимальный случай: однородный ближайший перенос на полном соответствии
Мориты $E=M_{20\times15}(\mathbb C)$ с двумя противоположными направлениями
распространения.

\section{Антикруговая сверка}

Этот тест не подставляет массы, матрицу смешивания или профиль дефекта.
Сначала требуется ответить на более дешёвый вопрос: фиксируют ли уже
имеющиеся бимодульная ковариантность, унитарность, локальность, чётность,
обращение времени и единый след сам однородный оператор переноса.

Литературная граница известна. Одномерный однородный причинный квантовый
автомат способен иметь уравнение Дирака в непрерывном пределе; см.
\path{arXiv:1212.2839}. Локальное квантовое блуждание может сходиться к
дираковской эволюции с контролируемой ошибкой; см.
\path{arXiv:1307.3524}. Но в этих конструкциях массовый параметр является
данным автомата. Поэтому проектный успех требует не воспроизвести известную
дисперсию, а вывести её смешивающий модуль из родительской архитектуры.

\section{Скалярность переноса на полном соответствии Мориты}

Рассмотрим общий прямоугольный бимодуль
\begin{equation}
 E_{p,q}=M_{p\times q}(\mathbb C)
\end{equation}
с естественными левым действием $M_p(\mathbb C)$ и правым действием
$M_q(\mathbb C)$. Пусть линейный оператор $T:E_{p,q}\to E_{p,q}$ является
бимодульным:
\begin{equation}
 T(a\xi b)=aT(\xi)b.
 \label{eq:v5-transfer-bimodule-map}
\end{equation}

\begin{lemma}[Скалярный коммутант полного прямоугольного модуля]
Для полных матричных углов
\begin{equation}
 \operatorname{End}_{M_p-M_q}(E_{p,q})\simeq\mathbb C.
 \label{eq:v5-full-morita-endomorphism}
\end{equation}
Следовательно, всякий $T$, удовлетворяющий
\eqref{eq:v5-transfer-bimodule-map}, имеет вид $T=\lambda I_E$.
\end{lemma}

Доказательство получается действием левых и правых матричных единиц:
значение $T$ на одной матричной единице определяет значения на всех
остальных, а совместимость с обоими диагональными действиями оставляет
только общий скаляр. Машинный аудит проверяет линейную систему на
представителе $M_{3\times2}(\mathbb C)$: матрица ограничений размера
$468\times36$ имеет ранг $35$, поэтому коммутант одномерен. Аргумент
матричных единиц не зависит от $p,q$ и переносится на $M_{20\times15}$.

Это одновременно сильный положительный и отрицательный результат.
Положительная часть состоит в том, что внутренний носитель размерности
$300$ не порождает сотни произвольных амплитуд переноса. Отрицательная
часть состоит в том, что полный бимодульный перенос сам по себе не различает
рёбра $u,d,e$.

\section{Двухнаправленный одношаговый оператор}

Для релятивистского переноса требуется как минимум пара встречных
хиральных направлений. После добавления ориентационного множителя
$\mathbb C^2$ наиболее общий ближайший однородный оператор, совместимый с
чётностью и обращением времени с точностью до несущественной общей фазы,
можно записать в импульсном представлении как
\begin{equation}
 W(k)=
 \begin{pmatrix}
  n e^{ik} & -im\\
  -im & n e^{-ik}
 \end{pmatrix}\otimes I_E,
 \qquad m,n\in\mathbb R.
 \label{eq:v5-minimal-dirac-walk}
\end{equation}
Чётность действует матрицей $\sigma_x$ и даёт
\begin{equation}
 \sigma_xW(k)\sigma_x=W(-k).
\end{equation}

Прямое умножение даёт
\begin{equation}
 W(k)W(k)^*=(m^2+n^2)I,
\end{equation}
поэтому унитарность эквивалентна единственному условию
\begin{equation}
 m^2+n^2=1.
 \label{eq:v5-transfer-unit-circle}
\end{equation}
После выбора ветви $n\ge0$ остаётся один непрерывный модуль
$m\in[-1,1]$.

Характеристическое уравнение унитарного оператора имеет вид
\begin{equation}
 \lambda^2-2n\cos k\,\lambda+1=0.
\end{equation}
Если $\lambda_\pm=e^{\pm i\omega}$, то точная дисперсия равна
\begin{equation}
 \cos\omega(k)=n\cos k.
 \label{eq:v5-transfer-dispersion}
\end{equation}

\begin{theorem}[Минимальный перенос с одним массовым модулем]
Полная ковариантность относительно $M_{20}$ и $M_{15}$ делает перенос
скалярным на $E=M_{20\times15}(\mathbb C)$. После минимального удвоения
направлений локальность, однородность, чётность, обращение времени и
унитарность оставляют семейство \eqref{eq:v5-minimal-dirac-walk} с одним
непрерывным параметром $m$. Текущие аксиомы не фиксируют его ненулевое
значение.
\end{theorem}

\section{Непрерывный дираковский предел}

Положим шаг решётки равным $a$ и выполним масштабирование
\begin{equation}
 m=aM,qquad k=ap,qquad n=\sqrt{1-a^2M^2}.
\end{equation}
Тогда
\begin{equation}
 n\cos k
 =1-\frac{a^2}{2}(M^2+p^2)+O(a^4).
\end{equation}
С другой стороны,
\begin{equation}
 \cos\!\left(a\sqrt{M^2+p^2}\right)
 =1-\frac{a^2}{2}(M^2+p^2)+O(a^4).
\end{equation}
Следовательно,
\begin{equation}
 \frac{\omega(ap)}a\longrightarrow\sqrt{M^2+p^2},
 \label{eq:v5-transfer-dirac-limit}
\end{equation}
и минимальная модель действительно имеет общий дираковский световой конус.
Машинная разность двух разложений до порядка $a^3$ тождественно равна
нулю.

Этот проход подтверждает пригодность языка переноса, но не выводит массу.
При $m=0$ получаются два независимых безмассовых хиральных сдвига. При
$|m|=1$ коэффициент пространственного переноса $n$ исчезает. Ненулевая
распространяющаяся масса требует
\begin{equation}
 0<|m|<1,
\end{equation}
то есть выбора точки на непрерывной окружности
\eqref{eq:v5-transfer-unit-circle}.

\section{Почему единый след не фиксирует массу}

Особенно привлекательна попытка использовать уже выведенные веса углов
связывающей алгебры:
\begin{equation}
 \tau_{35}(p_{20})=\frac47,
 \qquad
 \tau_{35}(p_{15})=\frac37.
 \label{eq:v5-transfer-corner-weights}
\end{equation}
Можно было бы предположить, например,
\begin{equation}
 m^2=\frac37,\quad n^2=\frac47.
 \label{eq:v5-tempting-weight-mass}
\end{equation}
Однако это не следует из сохранения следа. Сопряжение любым унитарным
$W(m)$ сохраняет нормированный след при каждом $m$. Более того, перестановка
ролей двух углов даёт столь же формально допустимое чтение
$m^2=4/7$, $n^2=3/7$.

Чтобы выбрать \eqref{eq:v5-tempting-weight-mass}, нужно дополнительно
объявить перенос отражением относительно очищения следового состояния,
стохастическим переходом или иной специальной функцией угловых весов. Ни
один из этих принципов пока не выведен из родителя Мориты.

\begin{nogocriterion}[Запрет подстановки углового веса в амплитуду]
Диагональные веса проекций $4/7$ и $3/7$ нельзя отождествлять с квадратами
внедиагональных амплитуд переноса без независимого функториального правила,
связывающего состояние на углах со связностью между углами.
\end{nogocriterion}

\section{Возвращение двух юкавских свобод}

Скалярный перенос на полном носителе имеет одну общую щель и потому не
создаёт иерархию между $u,d,e$. Чтобы различить три заряженных ребра, нужно
ограничить наблюдаемый угол и ввести
\begin{equation}
 m\longmapsto(m_u,m_d,m_e).
\end{equation}
Унитарность независимо определяет
\begin{equation}
 n_s=\sqrt{1-m_s^2},
 \qquad s\in\{u,d,e\},
\end{equation}
но не связывает три значения. После удаления общего масштаба остаются две
относительные координаты. Это в точности то же число свобод, которое было
найдено для правильно типизированного $H_{15}$-модуля одноформ.

\begin{proposition}[Динамическое возвращение статической неоднозначности]
Общий бимодульный перенос даёт одну универсальную щель и не воспроизводит
массовую иерархию. Блочное разделение переноса на $u,d,e$ вводит три
независимых массовых модуля и оставляет две относительные свободы после
общего масштаба. Поэтому минимальная динамическая модель не устраняет
прежнюю юкавскую неоднозначность, а воспроизводит её в дисперсионной форме.
\end{proposition}

\section{Почему дефект пока не вставляется}

Можно было бы теперь задать пространственно зависящий профиль $m(x)$,
поменять его знак и искать локализованную нулевую моду. Известные
дираковские и квантово-блуждательные модели допускают такую конструкцию.
Но профиль $m(x)$ уже содержал бы тот самый невыведенный модуль, который
требовалось получить. Его топологическая обработка показала бы следствия
заданного дефекта, но не происхождение дефекта.

Поэтому применяется ранняя остановка: гомогенный сектор уже выявил один
внешний непрерывный параметр. Нулевая мода, ароматные осцилляции и
майорановское ядро на этом шаге не объявляются результатами.

\section{Оставшаяся проектная подсказка}

В проекте имеется только один уже выведенный объект, который потенциально
может превратить фазовую информацию в каноническое действие на
ориентационной паре: ранг-один проектор вместе с тетраэдрической
голономией $2\pi/3$. Но одна голономная фаза фиксирует аргумент, а массовый
модуль $|m|$ является вещественной амплитудой. Поэтому следующий тест должен
быть предельно узким:

\begin{quote}
выводят ли проектор и тетраэдрическая голономия единственное отражение или
частичную изометрию на ориентационном дублете, модуль внедиагонального
элемента которой фиксирован без обращения к массам?
\end{quote}

Если такого функториального построения нет, минимальную динамическую ветвь
следует закрыть. Добавление произвольной монеты, профиля дефекта или
секторных амплитуд будет переименованной подгонкой.

\begin{protocol}
\begin{enumerate}
 \item бимодульный коммутант полного $M_{20\times15}$:
 \textbf{одномерен};
 \item локальный унитарный перенос:
 \textbf{построен};
 \item чётность и обращение направления:
 \textbf{выполнены};
 \item точная дисперсия:
 \textbf{$\cos\omega=n\cos k$};
 \item дираковский непрерывный предел:
 \textbf{получен};
 \item общий световой конус:
 \textbf{получен в масштабируемом пределе};
 \item единственный ненулевой массовый модуль:
 \textbf{не выведен};
 \item фиксация $m$ весами $4/7,3/7$:
 \textbf{не следует из следа};
 \item разделение $u,d,e$ без новых параметров:
 \textbf{не получено};
 \item топологический дефект и осцилляции:
 \textbf{не вводились по правилу ранней остановки};
 \item физическое замыкание:
 \textbf{не достигнуто}.
\end{enumerate}
\end{protocol}

Машинный сертификат сохранён в
\path{s2t_v5_local_defect_transfer_operator_gate.py} и
\path{s2t_v5_local_defect_transfer_operator_gate_results.json}.